\section{Related Work}
\label{sec:related}

Embedding compression splits into two families by \emph{when} the compression
happens. Post-hoc methods act on a frozen embedding. Training-time methods build
the compression into the objective. We review each, then position our study
against the four closest works.

\paragraph{Post-hoc compression.}
The simplest compressors transform an embedding after training, so they are cheap
and model-agnostic. Scalar quantization maps float32 to INT8 with per-dimension
calibration and is near-lossless at $4\times$~\citep{shakir2024embedding}. Binary
quantization keeps only the sign of each dimension and searches in Hamming space
at $32\times$. PQ~\citep{jegou2010product} splits the
vector into subspaces and encodes each with a learned codebook, and it remains
the strong learning-to-hash baseline. TurboQuant~\citep{zandieh2025turboquant}
is a recent data-oblivious quantizer that applies a random rotation before
scalar quantization to approach optimal distortion. The common limitation is
that none of these methods can change the representation they are handed.

\paragraph{Training-time compression.}
A second family reshapes the representation during training, at the cost of a
training run. Binarization is not differentiable, so it needs a gradient
surrogate. The STE~\citep{bengio2013estimating,
courbariaux2016bnn} passes gradients through the sign unchanged, while an annealed
tanh sharpens a smooth surrogate into a step function over
training~\citep{yamada2021efficient, yang2019quantization, schrodter2025trainable}.
MRL~\citep{kusupati2022matryoshka} takes a
different route and trains nested sub-vectors, so one model serves many truncation
dimensions. These methods can fix geometry that post-hoc methods cannot, but they
have mostly been studied one at a time.

\paragraph{The closest work falls in two camps.}
On the \emph{evaluation} side, CoRECT~\citep{correct2025} is nearest to us. It
compares compression techniques across corpus size and compression ratio for
several models, but only for post-hoc methods, so training-time compression is
absent. On the \emph{method} side, QAMA~\citep{qama2025} is nearest. It stacks
MRL with multi-bit quantization-aware training and learned thresholds, but it
relies on a penalty loss rather than a differentiable surrogate, and it never
compares gradient surrogates. Two more works touch single pieces of our study.
The Binary Passage Retriever (BPR)~\citep{yamada2021efficient} compares annealed
tanh against PQ and full precision, but never against STE. BEBR~\citep{gan2023binary}
deploys recurrent residual binarization at scale with runtime-selectable bit
depth, but trains the compressor separately from the encoder. Our own prior work
on a scientific-domain encoder~\citep{pantha2026indussde} shows that
binarization-aware training beats post-hoc binary, but on a single model, a single
domain, and STE alone.

\paragraph{Our position.}
No prior work compares post-hoc, training-time, \emph{and} stacked compression on
common ground. The broadest evaluation, CoRECT, varies the backbone but covers
only post-hoc methods; the training-time studies fix a single backbone or a
narrow pair. As a result, whether the best compression method depends on the
encoder's retrieval prior has not been tested, and the backbone-dependent
behavior we report has not been explained. We close that gap, and we connect the
behavior to the anisotropy of contextual
embeddings~\citep{ethayarajh2019contextual} through the effective
dimension~\citep{litwinkumar2017optimal, gao2017theory}.
